\chapter{Ablation training curves}
\label{app:curves}

For completeness, Figure~\ref{fig:ablation_curves} gives the training traces of
the three loss-reweighting arms of \S\ref{sec:ablation}, under the identical
quick-training configuration used for the ablation ($35$ epochs, leave-one-AR-out
on HARP~11930). The traces make the outcome of Table~\ref{tab:ablation} legible
epoch by epoch.

All three losses drive their training objective down quickly (left panel; the
delta loss plateaus higher only because its per-pixel weighting rescales the loss
magnitude, not because it fits worse). The held-out error (right panel) is where
they part company: the plain and horizon arms settle below the persistence
baseline, whereas the delta arm climbs \emph{above} it and stays there---the
model, told to ignore the flat regions that persistence already handles, degrades
them. This is the epoch-by-epoch view of the negative result reported in
\S\ref{sec:ablation}.

\begin{figure}[!htbp]
\centering
\includegraphics[width=\textwidth]{figures/fig_ablation_curves.pdf}
\caption[Ablation training curves]{\textbf{Training curves for the three loss
arms} (quick config, $35$ epochs, held-out HARP~11930). Left: Huber training
loss (log scale). Right: held-out mean absolute error against the persistence
baseline (dashed); the delta arm's epoch-$0$ value ($0.047$) is off-scale. Plain
and horizon converge below persistence; delta diverges above it.}
\label{fig:ablation_curves}
\end{figure}
